%TCIDATA{Version=5.50.0.2890}
%TCIDATA{LaTeXparent=0,0,NoticeDKE.tex}


The inversion procedure for solving the first order drift kinetic equation
is exactly similar to the one described in Sec. \ref{fpinversionsection}.\
Indeed, the linear system of equation to be solved may be cast in a similar
form. Usually, the convergence coefficient $\mathcal{R}_{\mathcal{C}}^{g}$
is set to similar values than for the Fokker-Planck problem, while the
maximum number of iterations allowed $\mathcal{N}_{\mathcal{C}}^{g}$ is
slightly larger \footnote{%
Usually twice $\mathcal{N}_{\mathcal{C}}^{f}$}, since the convergence
requires more iterations, around $20-30$. Since the function $g_{0}^{\left(
0\right) }$ may be negative and that $dt$ has no sense here, as the solution
is determined once the steady-state regime is achieved, the convergence
criterion is modified accordingly
\begin{equation}
\sqrt{\left| \frac{2\pi \iint \Delta g_{0}^{\left( 0\right) 2}g_{0}^{\left(
0\right) }J_{p}J_{\xi _{0}}dpd\xi _{0}}{2\pi \iint g_{0}^{\left( 0\right)
}J_{p}J_{\xi _{0}}dpd\xi _{0}}\right| }<\mathcal{R}_{\mathcal{C}}^{g}
\label{dkeconvcrit}
\end{equation}
at all plasma locations $\psi .$ Here again $J_{p}$ and $J_{\xi _{0}}$ are
the partial Jacobians introduced in Sec. \ref{fluxcons}\footnote{%
The possibility to define a fully $3-D$ criterion has been foreseen.
However, since collision rates may change drasticaly from plasma center to
the edge, this choice ensures that convergence is well achieved at the right
level at all plasma locations. }.

Projected on the numerical grid, the criterion (\ref{dkeconvcrit}) becomes

\begin{align}
&\left| \sum_{i=0}^{i=n_{p}-1}\sum_{j=0}^{j=n_{\xi _{0}}-1}\left[
g_{0,l+1/2,i+1/2,j+1/2}^{\left( 0\right) \left( k+1\right)
}-g_{0,l+1/2,i+1/2,j+1/2}^{\left( 0\right) \left( k\right)
}\right]
^{2}\right.\nonumber\\&\left.\quad\times g_{0,l+1/2,i+1/2,j+1/2}^{\left( 0\right) \left( k\right) }p_{i+1/2}^{2}%
\Delta p_{i+1/2}\Delta \xi
_{0,j+1/2}\right|^{1/2}\nonumber\\&\quad\times\left|
\sum_{i=0}^{i=n_{p}-1}%
\sum_{j=0}^{j=n_{\xi _{0}}-12}g_{0,l+1/2,i+1/2,j+1/2}^{\left( 0\right)
\left( k\right) }p_{i+1/2}^{2}\Delta p_{i+1/2}\Delta \xi _{0,j+1/2}\right|^{-1/2} %
<\mathcal{R}_{\mathcal{C}}^{g}
\end{align}
where $\left( k\right) $ is the iteration step number. Since all
quantities are normalized, this definition is fully acceptable. If
the number of iterations reaches $\mathcal{N}_{\mathcal{C}}^{g},$
the program warns that convergence is not achieved and that
something is going wrong in the calculations, as for the
Fokker-Planck part.
